% GA-04c.2c corrected derivative of cp_iv_0098_source_figure_01.tex
% Original source-backed figure retained unchanged for provenance.
% N=3: the positively oriented half-integer SQUARE has sides at +/- 7/2.
\begin{figure}[htbp]
  \centering
  \begin{tikzpicture}[scale=0.83]
    \draw[->] (-4.3,0) -- (4.3,0) node[below right] {$\Re z$};
    \draw[->] (0,-4.3) -- (0,4.3) node[above left] {$\Im z$};
    \draw[very thick,blue,->] (-3.5,-3.5) -- (3.5,-3.5);
    \draw[very thick,blue,->] (3.5,-3.5) -- (3.5,3.5);
    \draw[very thick,blue,->] (3.5,3.5) -- (-3.5,3.5);
    \draw[very thick,blue,->] (-3.5,3.5) -- (-3.5,-3.5);
    \foreach \n in {-3,-2,-1,0,1,2,3} {
      \fill[red] (\n,0) circle (2pt);
    }
    \fill[red] (0,1) circle (2pt);
    \fill[red] (0,-1) circle (2pt);
    \node[red,above right] at (0,1) {$i$};
    \node[red,below right] at (0,-1) {$-i$};
    \node[blue,above right] at (3.5,3.5) {$\Gamma_3$};
    \node[below] at (-3,0) {$-3$};
    \node[below] at (3,0) {$3$};
    \node[above left] at (-3.5,3.5) {$-\tfrac72+\tfrac72i$};
  \end{tikzpicture}
  \caption{The positively oriented half-integer square $\Gamma_3$ for the Fourier--residue calculation. All integer poles $-3,\dots,3$ and $\pm i$ lie strictly inside the contour.}
  \label{fig:cp-iv-0098-half-integer-square}
\end{figure}
